\begin{solapartat}
Com que la trompeta conté un extrem tancat i un altre obert, la condició per les possibles ones estacionàries dins de la seva cavitat és
\[ l_0 = \frac{\lambda_n}{4} + \frac{\lambda_n}{2}\,n = \frac{\lambda_n}{4}(2n+1) \quad n = 0, 1, 2, 3, \ldots \ \text{\punts{0,2}} \]
D'aquesta relació obtenim que les possibles ones estacionaries a la trompeta tenen longituds d'ona
\[ \lambda_n = \frac{4l_0}{2n+1} \ \text{\punts{0,2}} \]
Tanmateix, essent $\lambda = v/f$ on $v$ és la velocitat del so al medi i $f$ la freqüència de l'ona, resulta
\[ f_n = \frac{v}{\lambda_n} = \frac{v}{4l_0}(2n+1) \ \text{\punts{0,3}} \]
Això doncs, amb n=0, 1 i 2 obtenim els valors
\[ \left.\begin{aligned} &\qquad n = 0\\ \lambda_0 &= \tfrac{4\times 0,975}{1} = 3,90\ \mathrm{m}\\ f_0 &= \tfrac{340}{4\times 0,975} = 87,2\ \mathrm{Hz}\end{aligned}\right\} \text{\punts{0,1}}
\qquad
\left.\begin{aligned} &\qquad n = 1\\ \lambda_1 &= \tfrac{4\times 0,975}{3} = 1,30\ \mathrm{m}\\ f_1 &= \tfrac{340}{4\times 0,975}\,3 = 262\ \mathrm{Hz}\end{aligned}\right\} \text{\punts{0,1}} \]
\[ \left.\begin{aligned} &\qquad n = 2\\ \lambda_2 &= \tfrac{4\times 0,975}{5} = 0,78\ \mathrm{m}\\ f_2 &= \tfrac{340}{4\times 0,975}\,5 = 436\ \mathrm{Hz}\end{aligned}\right\} \text{\punts{0,1}} \]
\end{solapartat}

\begin{solapartat}
L'ona ressonant dins de la cavitat de la trompeta correspon al segon mode de vibració, es a dir, al mode $n = 1$ \punts{0,2} de les expressions anteriors. Això doncs hauria de ser $l = 3\lambda/4$. Com que $\lambda = v/f$, resulta
\[ l = \frac{3}{4}\lambda = \frac{3}{4}\left(\frac{v}{f}\right) = \frac{3\times 340}{4\times 247} = 1,03\ \mathrm{m}\,. \ \text{\punts{0,4}} \]
La variació en la longitud de la cavitat recorreguda per l'aire quan es prem el segon pistó és, per tant,
\[ l_1 = l - l_0 = 1,03 - 0,975 = 5,5\cdot 10^{-2}\ \mathrm{m} \ \text{\punts{0,4}} \]
\end{solapartat}
